\subsection{Manipulación de cadenas}

Las cadenas (\texttt{str}) son inmutables en Python, por lo que todas las operaciones que modifican su contenido devuelven una nueva cadena.

La mayoría de problemas de parsing, simulación y procesamiento de texto utilizan estas funciones.

\vspace{0.3cm}

\subsubsection{replace()}

Reemplaza ocurrencias de una subcadena.

\begin{lstlisting}[language=Python]
s = "banana"
s.replace("a", "o")
# "bonono"
s.replace("a", "o", 2)
# "bonona"
\end{lstlisting}

Complejidad promedio: \bigO{n}

\vspace{0.3cm}

\subsubsection{split()}

Divide una cadena utilizando un separador.

\begin{lstlisting}[language=Python]
s.split()
s.split(",")
s.split(":")
s.split("\n")
\end{lstlisting}

Si no se especifica un separador, utiliza cualquier cantidad de espacios consecutivos.

\begin{lstlisting}[language=Python]
s = "a    b     c"
print(s.split())
# ['a', 'b', 'c']
\end{lstlisting}

\vspace{0.3cm}

\subsubsection{join()}

Une una colección de cadenas utilizando un separador.

Es mucho más eficiente que concatenar strings dentro de un ciclo.

\begin{lstlisting}[language=Python]
nums = ["1", "2", "3"]
",".join(nums)
# "1,2,3"
" ".join(nums)
# "1 2 3"
\end{lstlisting}

\graybold{Complejidad:} \bigO{n}

\vspace{0.3cm}

\subsubsection{strip(), lstrip() y rstrip()}

Eliminan caracteres al inicio y/o final.

\begin{lstlisting}[language=Python]
s.strip()
s.lstrip()
s.rstrip()
\end{lstlisting}

También pueden eliminar caracteres específicos.

\begin{lstlisting}[language=Python]
"--ICPC--".strip("-")
# ICPC
\end{lstlisting}

\vspace{0.3cm}

\subsubsection{find() y rfind()}

Buscan una subcadena.

Devuelven

\begin{itemize}
\item la posición encontrada;
\item -1 si no existe.
\end{itemize}

\begin{lstlisting}[language=Python]
s.find("abc")
s.rfind("abc")
\end{lstlisting}

\vspace{0.3cm}

\subsubsection{count()}

Cuenta cuántas veces aparece una subcadena.

\begin{lstlisting}[language=Python]
s.count("ab")
\end{lstlisting}

\vspace{0.3cm}

\subsubsection{startswith() y endswith()}

Permiten verificar prefijos y sufijos.

\begin{lstlisting}[language=Python]
s.startswith("abc")
s.endswith(".cpp")
\end{lstlisting}

Muy útiles para parsing.

\vspace{0.3cm}

\subsubsection{lower(), upper() y capitalize()}

Cambian el formato de las letras.

\begin{lstlisting}[language=Python]
s.lower()
s.upper()
s.capitalize()
\end{lstlisting}

\vspace{0.3cm}

\subsubsection{isalpha(), isdigit(), isalnum()}

Verifican el contenido de una cadena.

\begin{lstlisting}[language=Python]
s.isalpha()
s.isdigit()
s.isalnum()
s.isspace()
\end{lstlisting}

\vspace{0.3cm}

\subsubsection{partition()}

Divide únicamente en la primera aparición del separador.

\begin{lstlisting}[language=Python]
s.partition("=")
# ('x', '=', '15')
\end{lstlisting}

A diferencia de split(), siempre devuelve exactamente tres elementos.

\vspace{0.3cm}

\subsubsection{Concatenación}

\begin{lstlisting}[language=Python]
a + b
a += b
\end{lstlisting}

Cuando se concatenan muchas cadenas es preferible utilizar

\begin{lstlisting}[language=Python]
"".join(lista)
\end{lstlisting}

ya que evita crear múltiples copias intermedias.

\vspace{0.3cm}

\subsubsection{Formateo de cadenas}

Python permite insertar variables directamente mediante f-strings.

\begin{lstlisting}[language=Python]
x = 3.14159265
f"{x:.2f}"
# 3.14
f"{x:.10f}"
# 3.1415926500
nombre = "ICPC"
f"Hola {nombre}"
\end{lstlisting}

Este mecanismo es el recomendado para imprimir números reales.

\vspace{0.3cm}

\subsubsection{Aplicaciones frecuentes en ICPC}

Eliminar espacios al leer una línea.

\begin{lstlisting}[language=Python]
s = input().strip()
\end{lstlisting}

Separar una línea.

\begin{lstlisting}[language=Python]
tokens = input().split()
\end{lstlisting}

Unir la respuesta.

\begin{lstlisting}[language=Python]
print(" ".join(ans))
\end{lstlisting}

Comprobar un prefijo.

\begin{lstlisting}[language=Python]
if s.startswith("YES"):
    ...
\end{lstlisting}

Eliminar caracteres.

\begin{lstlisting}[language=Python]
s = s.replace(",", "")
\end{lstlisting}

Buscar una palabra.

\begin{lstlisting}[language=Python]
if s.find("abc") != -1:
    ...
\end{lstlisting}

Contar caracteres.

\begin{lstlisting}[language=Python]
vocales = s.count("a")
\end{lstlisting}

\vspace{0.3cm}

\subsubsection*{Resumen}

\begin{center}
\begin{tabular}{|l|l|}
\hline
Función & Uso principal \\ \hline
replace() & Reemplazar texto \\ \hline
split() & Dividir una cadena \\ \hline
join() & Unir cadenas \\ \hline
strip() & Eliminar espacios \\ \hline
find() & Buscar subcadena \\ \hline
rfind() & Buscar desde el final \\ \hline
count() & Contar ocurrencias \\ \hline
startswith() & Verificar prefijo \\ \hline
endswith() & Verificar sufijo \\ \hline
lower() & Minúsculas \\ \hline
upper() & Mayúsculas \\ \hline
capitalize() & Primera letra mayúscula \\ \hline
isalpha() & Solo letras \\ \hline
isdigit() & Solo dígitos \\ \hline
isalnum() & Alfanumérico \\ \hline
isspace() & Solo espacios \\ \hline
partition() & Dividir una vez \\ \hline
\end{tabular}
\end{center}